\documentclass{ximera}

\input{../../preamble.tex}

\title{Chaining Functions}

\begin{document}

\begin{abstract}
Chain Rule
\end{abstract}
\maketitle





Using two given two functions, we can create a third function, called the \textbf{composition}.

If $f$ and $g$ are the names of the two component functions, then $f \circ g$ is the symbol for the composition of $f$ and $g$.

We compose $f$ and $g$ to form $(f \circ g)$.

The order makes a difference.

We compose $g$ and $f$ to form $(g \circ f)$.


Most of time these are not equal functions.

\[
(f \circ g) \ne (g \circ f)
\]


\begin{idea} \textbf{\textcolor{blue!55!black}{Composition}}


Let $F$ and $G$ be two functions.

The composition of $F$ and $G$ is another function built from pieces of $F$ and $G$.  A circle between the function names denotes a composition.

\[
F \circ G
\]


\[
(F \circ G)(a) = F(G(a))
\]


The composition applies $G$ to an element of the domain of $G$ and then uses the resulting function value from $G$ as a domain member of $F$ obtaining a function value of $F$.

This ultimate value of $F$ resulting from this process is the value of $F \circ G$.

\end{idea}
The domain of $F \circ G$ is inside the domain of $G$.\\
The range of $F \circ G$ is inside the range of $F$.


Composition is our way of representing an ordered application of separate processes.



\subsection*{Graphical Composition}


When representing a function graphically, the domain sits on the horizontal axis and the range sits on the vertical axis.

The composing process begins with a number on the inner function's horizontal axis. The inner function's value is found on the vertical axis via the points on the graph.

Attention then turns to the horizontal axis of the graph for the outer function. Locate the value of the inner function on the domain axis of the outer function.

Now, run through the evaluation process once more with the outer function.

From this outer function domain number, the value of the outer function is determined via the points on the graph of the outer function.

This final number on the vertical axis of the outer function is the value of the composition.



\textbf{Two Component Functions:}


Let $f$ and $w$ be two functions defined by the following graphs.



Below is the graph of $y=f(x)$.  

\begin{image}
\begin{tikzpicture} 
  \begin{axis}[
            domain=-10:10, ymax=10, xmax=10, ymin=-10, xmin=-10,
            axis lines =center, xlabel=$x$, ylabel={$y=f(x)$},
            grid = both, 
            ytick={-10,-8,-6,-4,-2,2,4,6,8,10}, 
            xtick={-10,-8,-6,-4,-2,2,4,6,8,10},
            yticklabels={$-10$,$-8$,$-6$,$-4$,$-2$,$2$,$4$,$6$,$8$,$10$}, 
            xticklabels={$-10$,$-8$,$-6$,$-4$,$-2$,$2$,$4$,$6$,$8$,$10$},
            ticklabel style={font=\scriptsize},
            every axis y label/.style={at=(current axis.above origin),anchor=south},
            every axis x label/.style={at=(current axis.right of origin),anchor=west},
            axis on top
          ]
          
          \addplot [line width=2, penColor, smooth,samples=100,domain=(-5:0)] ({x},{(x+5)*(x-1)});
          \addplot [line width=2, penColor, smooth,samples=100,domain=(1:5)] ({x},{x-3});
          \addplot [line width=2, penColor, smooth,samples=100,domain=(5:8)] ({x},{2*x-16});

          \addplot [color=penColor,fill=white,only marks,mark=*] coordinates{(0,-5) (1,-2) (5, 2) (8,0)};
          \addplot [color=penColor,only marks,mark=*] coordinates{(-5,0) (5,-6)};


           

  \end{axis}
\end{tikzpicture}
\end{image}






\textbf{\textcolor{blue!55!black}{$\blacktriangleright$ desmos graph}} 
\begin{center}
\desmos{1hdvwizmhq}{400}{300}
\end{center}









Below is the graph of $y=w(t)$. 

\begin{image}
\begin{tikzpicture} 
  \begin{axis}[
            domain=-10:10, ymax=10, xmax=10, ymin=-10, xmin=-10,
            axis lines =center, xlabel=$t$, ylabel={$y=w(t)$},
            grid = both, 
            ytick={-10,-8,-6,-4,-2,2,4,6,8,10}, 
            xtick={-10,-8,-6,-4,-2,2,4,6,8,10},
            yticklabels={$-10$,$-8$,$-6$,$-4$,$-2$,$2$,$4$,$6$,$8$,$10$}, 
            xticklabels={$-10$,$-8$,$-6$,$-4$,$-2$,$2$,$4$,$6$,$8$,$10$},
            ticklabel style={font=\scriptsize},
            every axis y label/.style={at=(current axis.above origin),anchor=south},
            every axis x label/.style={at=(current axis.right of origin),anchor=west},
            axis on top
          ]
          
          \addplot [line width=2, penColor, smooth,samples=100,domain=(-7:-2)] ({x},{-x-1});
          \addplot [line width=2, penColor, smooth,samples=100,domain=(-2:6)] ({x},{-0.5*x-1});
          \addplot [line width=2, penColor, smooth,samples=100,domain=(6:8)] ({x},{-2*x+7});

          \addplot [color=penColor,fill=white,only marks,mark=*] coordinates{(-7,6) (-2,1) (6,-4)};
          \addplot [color=penColor,only marks,mark=*] coordinates{(-2,0) (6,-5) (8,-9)};


           

  \end{axis}
\end{tikzpicture}
\end{image}







\textbf{\textcolor{blue!55!black}{$\blacktriangleright$ desmos graph}} 
\begin{center}
\desmos{pr7lf81kkh}{400}{300}
\end{center}






\begin{example}

From the graph of $f$, we can see that $f(-4) = -5$.

From the graph of $w$, we can see that $w(-5) = 4$.

These two values tell us that $(w \circ f)(-4) = 4$

\end{example}






\begin{example}

From the graph of $f$, we can see that $f(7) = -2$.

From the graph of $w$, we can see that $w(-2) = 0$.

These two values tell us that $(w \circ f)(7) = 0$

\end{example}




\begin{example}

From the graph of $w$, we can see that $w(2) = -2$.

From the graph of $f$, we can see that $f(-2) = -9$.

These two values tell us that $(f \circ w)(2) = -9$

\end{example}





\begin{example}

From the graph of $w$, we can see that $w(-6) = 5$.

From the graph of $f$, we can see that $f(5) = -6$.

These two values tell us that $(f \circ w)(-6) = -6$

\end{example}








\begin{example}

From the graph of $w$, we can see that $w(-2) = 0$.

From the graph of $f$, we can see that $f(0)$ is undefined.

Therefore, the domain of $(f \circ w)$ cannot include $-2$.

\end{example}





\begin{question}

From the graph of $f$, we can see that $f(3) = \answer[tolerance=0.1]{0}$.

From the graph of $w$, we can see that $w(0) = \answer[tolerance=0.1]{-1}$.

These two values tell us that $(w \circ f)(3) = \answer[tolerance=0.1]{-1}$

\end{question}







\begin{question}

From the graph of $w$, we can see that $w(4) = \answer[tolerance=0.1]{-3}$.

From the graph of $f$, we can see that $f(-3) = \answer[tolerance=0.1]{-8}$.

These two values tell us that $(f \circ w)(4) = \answer[tolerance=0.1]{-8}$

\end{question}

























\subsection*{The Chain Rule}


The Chain Rule tells us when composition functions are increasing and decreasing.

The Chain Rule combines the behavior of the component functions and identifies the behavior of the composite function.

This allows us to study very complex functions by examining the pieces from which they are built.




\begin{fact}  \textbf{\textcolor{red!75!green}{$inc \circ inc = inc$}} 

Let $f$ be an increasing function. \\
Let $g$ be an increasing function.


Let $a < b$.

Since $f$ is an increasing function, we know that $f(a) < f(b)$.

Now evaluate $g$ at $f(a)$ and $f(b)$.


We have $f(a) < f(b)$ and $g$ is an increasing function. 

That give us $g(f(a)) < g(f(b))$



\[ \text{If } \, a < b \, \text{ then, } \, g(f(a)) < g(f(b)) \]

That tells us that $g \circ f$ is an increasing function.

\end{fact}








\begin{fact}  \textbf{\textcolor{red!75!green}{$dec \circ inc = dec$}} 

Let $f$ be an increasing function. \\
Let $g$ be a decreasing function.


Let $a < b$.

Since $f$ is an increasing function, we know that $f(a) < f(b)$.

Now evaluate $g$ at $f(a)$ and $f(b)$.


We have $f(a) < f(b)$ and $g$ is a decreasing function. 

That give us $g(f(a)) > g(f(b))$



\[ \text{If } \, a < b \, \text{ then, } \, g(f(a)) > g(f(b)) \]

That tells us that $g \circ f$ is a decreasing function.

\end{fact}















\begin{fact}  \textbf{\textcolor{red!75!green}{$inc \circ dec = dec$}} 

Let $f$ be a decreasing function. \\
Let $g$ be an increasing function.


Let $a < b$.

Since $f$ is a decreasing function, we know that $f(a) > f(b)$.

Now evaluate $g$ at $f(a)$ and $f(b)$.


We have $f(a) > f(b)$ and $g$ is an increasing function. 

That give us $g(f(a)) > g(f(b))$



\[ \text{If } \, a < b \, \text{ then, } \, g(f(a)) > g(f(b)) \]

That tells us that $g \circ f$ is a decreasing function.

\end{fact}
















\begin{fact}  \textbf{\textcolor{red!75!green}{$dec \circ dec = inc$}} 

Let $f$ be a decreasing function. \\
Let $g$ be a decreasing function.


Let $a < b$.

Since $f$ is a decreasing function, we know that $f(a) > f(b)$.

Now evaluate $g$ at $f(a)$ and $f(b)$.


We have $f(a) > f(b)$ and $g$ is a decreasing function. 

That give us $g(f(a)) < g(f(b))$



\[ \text{If } \, a < b \, \text{ then, } \, g(f(a)) < g(f(b)) \]

That tells us that $g \circ f$ is an increasing function.

\end{fact}


These are the four combinations that make up \textbf{\textcolor{-red!75!green}{The Chain Rule}}.


\begin{theorem} \textbf{\textcolor{-red!75!green}{The Chain Rule}} 

\[ inc \circ inc = inc \]

\[ dec \circ inc = dec \]

\[ inc \circ dec = dec \]

\[ dec \circ dec = inc \]

\end{theorem}



The Chain Rule is indispensable.

Without it, we could never analyze any meaningful functions.

We will slowly uncover its secrets and how to apply it in any situation.


















\begin{onlineOnly}
\begin{center}
\textbf{\textcolor{green!50!black}{ooooo-=-=-=-ooOoo-=-=-=-ooooo}} \\

more examples can be found by following this link\\ \link[More Examples of Formulas]{https://ximera.osu.edu/csccmathematics/precalculus1/precalculus1/formulas/examples/exampleList}

\end{center}
\end{onlineOnly}





\end{document}
